\documentclass[11pt, answers]{exam}
\usepackage[empty]{fullpage}
\usepackage{amsmath,amssymb}
\usepackage{colortbl}
\usepackage{subfig}
\usepackage{multicol}
\usepackage{tikz,pgfplots}
\usepgfplotslibrary{statistics}
\usepackage{color,comment,soul}

\newcommand{\blank}[1]{\underline{\hspace*{#1}}}
\newcommand{\ds}{\displaystyle}
\newcommand{\inner}[1]{\langle {#1} \rangle}
\newcommand{\R}{\mathbb{R}}
\newcommand{\on}{\operatorname}


\begin{document}
\graphicspath{{/home/brian/Dropbox/HSC/Spring16/Math111/}}

\subsection*{Numerical Analysis - Math 342 \hfill Midterm 2 Review}

\textit{The following problems are similar to ones you might see on the midterm exam.}

\begin{questions}

\question How many floating point operations does it take to compute $y = a + bx + cx^2$ if $a, b, c$ and $x$ are all floating point numbers? 
\begin{solution}
There are two additions, and three multiplications (including one to compute $x^2$) for a total of 5 flops. 
\end{solution}
\vfill

\question It is possible to re-write the polynomial above as $y = a + x(b + cx)$.  How many floating point operations would it take to compute this alternative formula for the polynomial? 
\begin{solution}
Now there are two additions and two multiplications for a total of 4 flops. 
\end{solution}
\vfill



\question Use the method of divided differences to find the Newton basis interpolating polynomial for the points $(0,0)$, $(1,1)$, and $(4,2)$.   

\begin{solution}
The table of divided differences is 

\begin{tabular}{c|ccc}
x & y & DD1 & DD2 \\ \hline
0 & 0 & ~  & ~ \\
~ & ~ & 1  & ~ \\
1 & 1 & ~  & $-\tfrac{1}{6}$ \\
~ & ~ & $\tfrac{1}{3}$  & ~ \\
4 & 2 & ~  & ~ \\
\end{tabular}

\noindent
So the interpolating polynomial is $p_2(x) = x - \tfrac{1}{6}x(x-1)$ in the Newton basis.
\end{solution}
\vfill

\question What is the interpolating polynomial above written in terms of the Lagrange basis polynomials? 
\begin{solution}
The Lagrange basis polynomials are
$$L_0(x) = \tfrac{1}{4}(x-1)(x-4) ~~~~ L_1(x) = -\tfrac{1}{3} x(x-4) ~~~~ L_2(x) = \tfrac{1}{12} x(x-1)$$
so the interpolating polynomial is $p_2(x) = -\tfrac{1}{3} x(x-4) + \tfrac{1}{6}x(x-1)$ in the Lagrange basis.
\end{solution}
\vfill

\newpage
\question Write down the Vandermonde matrix system for these same points $(0,0)$, $(1,1)$, $(4,2)$ to find the coefficients of the interpolating polynomial in the standard basis. You don't need to solve the Vandermonde matrix system.
\begin{solution}
$$\begin{bmatrix} 1 & 0 & 0 \\ 1 & 1 & 1 \\ 1 & 4 & 16\end{bmatrix} \begin{bmatrix} c_0 \\ c_1 \\ c_2 \end{bmatrix} = \begin{bmatrix} 0 \\ 1 \\ 2 \end{bmatrix}.$$
\end{solution}
\vfill


\question The second degree interpolating polynomial for the function $f(x) = 1/x$ on the interval $[1,3]$ with three equally spaced nodes ($x_0 = 1$, $x_1=2$, $x_2 =3$) is 
$$p_2(x) = \tfrac{1}{6}x^2 - x + \tfrac{11}{6}.$$
The error formula for interpolating polynomials with equally spaced nodes is 
$$|f(x)-p_n(x)| \le \frac{h^{n+1}}{4(n+1)} \max_{\xi \in [a,b]} |f^{(n+1)}(\xi)|$$ 
where $h = \tfrac{b-a}{n}$.  Use this formula to find an upper bound on the error in using the polynomial $p_2$ to approximate $f(x) = 1/x$ on the interval $[1,3]$.  
\begin{solution}
In this example $n = 2$ and $h = 1$.  The third derivative of $f$ is $f^{(3)}(x) = -6 x^{-4}$. The largest possible absolute value for this function would occur at the left endpoint when $x = 1$, so the error is bounded by:
$$ \frac{h^{n+1}}{4(n+1)} \max_{\xi \in [a,b]} |f^{(n+1)}(\xi)| \le \frac{1}{12} 6(1^{-4}) = \frac{1}{2}.$$ 
\end{solution}
\vfill



\question Let $\mathbf{x} \in \R^n$ be any vector with $\|\mathbf{x}\| = \sqrt{\mathbf{x}^T \mathbf{x}} = 1$.  Let $I$ denote the $n$-by-$n$ identity matrix.  
\begin{parts} 
\part Simplify the following expression as much as possible. 
$$(I - 2\mathbf{x}\mathbf{x}^T) (I - 2 \mathbf{x}\mathbf{x}^T).$$
\begin{solution}
First FOIL
$$(I - 2\mathbf{x}\mathbf{x}^T) (I - 2 \mathbf{x}\mathbf{x}^T) = I - 2 \mathbf{x}\mathbf{x}^T - 2 \mathbf{x}\mathbf{x}^T + 4 \mathbf{x}\mathbf{x}^T \mathbf{x} \mathbf{x}^T$$ 
Then cancel $\mathbf{x}^T \mathbf{x}$ in the last term, since that is just 1. 
$$=  I - 2 \mathbf{x}\mathbf{x}^T - 2 \mathbf{x}\mathbf{x}^T + 4 \mathbf{x}\mathbf{x}^T $$
Then collect the like terms and everything cancels except for $I$. So the answer is $I$.
\end{solution}
\vfill 

\part What does the answer to part (a) imply about the columns of the matrix $I - 2 \mathbf{x} \mathbf{x}^T$? 
\begin{solution}
It means that the columns of $I - 2 \mathbf{x} \mathbf{x}^T$ are an orthonormal set. 
\end{solution}
\vfill
\end{parts}


\newpage
\question The normal equation to find the coefficients of a (discrete) least square regression model is 
$$X^T X b = X^T y.$$
Suppose you want the best fit linear function $\hat{y} = b_0 + b_1 x$ to approximate the points $(-2,3)$, $(0, 2)$, $(2,0)$.  
\begin{parts}
\part What is the matrix $X$ and the vector $y$ in the normal equation above?

\begin{solution}
$$X = \begin{bmatrix} 1 & -2 \\ 1 & 0 \\ 1 & 2 \end{bmatrix} ~~~ \text{ and } ~~~ y = \begin{bmatrix} 3 \\ 2\\ 0 \end{bmatrix}.$$
\end{solution}
\vfill

\part Compute $X^TX$ and $X^T y$.  
\begin{solution}
$$X^TX = \begin{bmatrix} 3 &  0 \\ 0 & 8  \end{bmatrix} ~~~ \text{ and } ~~~ X^Ty = \begin{bmatrix} 5 \\ -6 \end{bmatrix}.$$
\end{solution}
\vfill

\part Solve the normal equations to find the coefficients of the regression line $\hat{y} = b_0 + b_1 x.$

\begin{solution}
The slope is $b_1 = -3/4$ and the y-intercept is $b_0 = 5/3$.  
\end{solution}
\vfill

\end{parts}

\question The Legendre polynomials are a family of orthogonal functions on the interval $[-1,1]$.  The first three Legendre polynomials are 
$$P_0(x) = 1 ~~~~ P_1(x) = x ~~~~ P_2(x) = x^2 - \tfrac{1}{3}.$$
Using the Legendre polynomials as a basis, find the best fit (continuous least squares) 2nd degree polynomial approximation of the function $\cos x$ on the interval $[-1,1]$. Express your answer with coefficients that are accurate to 4 significant digits. You can use the following integrals to help
$$\int_{-1}^1 P_0(x) \cos x \,dx = 1.683 ~~~~ \int_{-1}^1 P_1(x) \cos x \,dx = 0$$
and
$$\int_{-1}^1 P_2(x) \cos x \,dx = -0.08271 ~~~~ \int_{-1}^1 P_2(x)^2 \, \, dx = \frac{8}{45}.$$ 
\begin{solution}
The least squares solution is 
$$\frac{\inner{P_0,\cos x}}{\inner{P_0,P_0}} P_0(x) + \frac{\inner{P_1,\cos x}}{\inner{P_1,P_1}} P_1(x) +\frac{\inner{P_2,\cos x}}{\inner{P_2,P_2}} P_2(x)$$
Four of those inner-products were given as integrals above.  The only two left to calculate are:
$$\inner{P_0,P_0} = \int_{-1}^1 1 \,dx = 2$$
and
$$\inner{P_1,P_1} = \int_{-1}^1 x^2 \,dx = \frac{2}{3}.$$
So the final answer is 
$$\frac{1.683}{2} P_0(x) - \frac{0.08271}{8/45} P_2(x) = 0.8415 - 0.4652 (x^2- \tfrac{1}{3}).$$
\end{solution}
\vfill

\newpage

\question Apply the Gram-Schmidt algorithm to the vectors to get an orthogonal basis for $\R^3$. 
$$\mathbf{x}_1 = \begin{bmatrix} 1 \\ 1 \\ 0 \end{bmatrix}, \mathbf{x}_2 = \begin{bmatrix} 3 \\ -1 \\ 1 \end{bmatrix}, \mathbf{x}_3 = \begin{bmatrix} 0 \\ 2 \\ -5 \end{bmatrix}.$$
\begin{solution}
$$\mathbf{v}_1 = \mathbf{x}_1 = \begin{bmatrix} 1 \\ 1 \\ 0 \end{bmatrix}$$
$$\mathbf{v}_2 = \mathbf{x}_2 - \frac{\inner{\mathbf{v_1}, \mathbf{x}_2}}{\inner{\mathbf{v}_1, \mathbf{v}_1}} \mathbf{v}_1 = \begin{bmatrix} 3 \\ -1 \\ 1 \end{bmatrix} - \frac{2}{2} \begin{bmatrix} 1 \\ 1 \\ 0 \end{bmatrix} = \begin{bmatrix} 2 \\ -2 \\ 1 \end{bmatrix}$$
$$\mathbf{v}_3 = \mathbf{x}_3 - \frac{\inner{\mathbf{v_1}, \mathbf{x}_3}}{\inner{\mathbf{v}_1, \mathbf{v}_1}} \mathbf{v}_1 -  \frac{\inner{\mathbf{v_2}, \mathbf{x}_3}}{\inner{\mathbf{v}_2, \mathbf{v}_2}} \mathbf{v}_2 = \begin{bmatrix} 0 \\ 2 \\ -5 \end{bmatrix} - \frac{2}{2} \begin{bmatrix} 1 \\ 1 \\ 0 \end{bmatrix}  - \frac{-9}{9} \begin{bmatrix} 2 \\ -2 \\ 1 \end{bmatrix} = \begin{bmatrix} 1 \\ -1 \\ -4 \end{bmatrix}.$$
\end{solution}
\vfill
\vfill

\question Let $\mathbf{x} = \begin{bmatrix} 1 \\ 1 \\ 1\\ 1 \end{bmatrix}$ and let $\mathbf{y} = \begin{bmatrix} 5 \\ 3 \\ 1 \\ -1 \end{bmatrix}$. 
\begin{parts}
\part Find the orthogonal projection of $\mathbf{y}$ onto the span of $\mathbf{x}$. 
\begin{solution}
$$\on{Proj}_\mathbf{x}(\mathbf{y}) = \frac{\inner{\mathbf{y}, \mathbf{x}}}{\inner{\mathbf{x}, \mathbf{x}}} \mathbf{x} = \frac{8}{4} \begin{bmatrix} 1 \\ 1 \\ 1 \\ 1 \end{bmatrix} = \begin{bmatrix} 2 \\ 2 \\ 2\\ 2 \end{bmatrix}.$$
\end{solution}
\vfill

\part Apply the Gram-Schmidt algorithm to $\{ \mathbf{x}, \mathbf{y} \}$ to get two orthogonal vectors with the same span. 
\begin{solution}
$$\mathbf{v}_1 = \mathbf{x} = \begin{bmatrix} 1 \\ 1\\ 1 \\ 1 \end{bmatrix}$$
and 
$$\mathbf{v}_2 = \mathbf{y} - \on{Proj}_\mathbf{x}(\mathbf{y}) = \begin{bmatrix} 5 \\ 3\\ 1 \\ -1 \end{bmatrix} - \begin{bmatrix} 2 \\ 2\\ 2\\ 2 \end{bmatrix} = \begin{bmatrix} 3 \\ 1 \\ -1 \\ -3 \end{bmatrix}.$$
\end{solution}
\vfill
\end{parts}

\question Find the inner product of $f(x) = 3x^2 - 1$ and $g(x) = x^2$ on the interval $[-1,1]$. 
\begin{solution}
$$\inner{f, g} = \int_{-1}^1 3 x^4 - x^2 \,dx = 2 \left[ \frac{3x^5}{5} - \frac{x^3}{3} \right]_{0}^1 = \frac{3}{5} - \frac{1}{3} = \frac{4}{15}.$$
\end{solution}
\vfill

\question Explain how you can tell that the functions $f(x) = x^4 + 5$ and $g(x) = \sin x$ are orthogonal in $L^2[-1,1]$ without actually integrating.  
\begin{solution}
Because $f(x)$ is an even function and $g(x)$ is an odd function, the product $f(x) g(x)$ is an odd function. So its integral from $-1$ to 1 will be zero. 
\end{solution}
\vfill 
\end{questions}

\end{document}
